\documentclass[12pt,a4paper]{article}
\usepackage[a4paper,left=25mm,right=20mm,top=24mm,bottom=24mm]{geometry}
\usepackage{fontspec}
\usepackage{polyglossia}
\setdefaultlanguage{russian}
\setotherlanguage{english}
\IfFontExistsTF{Libertinus Serif}{
    \setmainfont{Libertinus Serif}
    \newfontfamily\cyrillicfont{Libertinus Serif}
}{
    \IfFontExistsTF{DejaVu Serif}{
        \setmainfont{DejaVu Serif}
        \newfontfamily\cyrillicfont{DejaVu Serif}
    }{
        \setmainfont{Times New Roman}
        \newfontfamily\cyrillicfont{Times New Roman}
    }
}
\IfFontExistsTF{DejaVu Sans Mono}{\setmonofont{DejaVu Sans Mono}[Scale=0.82]}{}
\usepackage{microtype}
\usepackage{setspace}
\setstretch{1.20}
\usepackage{parskip}
\setlength{\parindent}{0pt}
\usepackage{booktabs}
\usepackage{tabularx}
\usepackage{longtable}
\usepackage{array}
\usepackage{ragged2e}
\newcolumntype{P}[1]{>{\RaggedRight\arraybackslash}p{#1}}
\newcolumntype{Y}{>{\RaggedRight\arraybackslash}X}
\setlength{\emergencystretch}{3em}
\usepackage{enumitem}
\usepackage{needspace}
\usepackage{amsmath}
\usepackage{amssymb}
\setlist{topsep=3pt,itemsep=2pt,parsep=0pt}
\usepackage{graphicx}
\usepackage{xcolor}
\definecolor{codegray}{RGB}{62,66,72}
\usepackage{listings}
\lstset{
    basicstyle=\ttfamily\footnotesize,
    breaklines=true,
    breakatwhitespace=false,
    columns=fullflexible,
    showstringspaces=false,
    frame=single,
    rulecolor=\color{black!20},
    xleftmargin=4pt,
    xrightmargin=4pt,
    aboveskip=8pt,
    belowskip=8pt
}
\usepackage{tikz}
\usetikzlibrary{arrows.meta,positioning,shapes.geometric}
\usepackage{hyperref}
\hypersetup{hidelinks,pdfauthor={Agibalov Vladislav},pdflang={ru-RU}}
\usepackage{fancyhdr}
\pagestyle{fancy}
\setlength{\headheight}{14pt}
\fancyhf{}
\fancyhead[L]{\small mangen}
\fancyhead[R]{\small Техническая документация · 2026}
\fancyfoot[C]{\thepage}
\renewcommand{\headrulewidth}{0.2pt}
\setcounter{tocdepth}{2}
\newcommand{\impl}{\textbf{Реализация.}~}

\title{mangen\\Архитектура и внутренние контракты}
\author{Агибалов Владислав Сергеевич\\НИЯУ МИФИ, ИВТ}
\date{2026 год}

\begin{document}
\maketitle
\begin{abstract}
Архитектура UNIX-утилиты \texttt{mangen} на C17 и POSIX API: состав модулей,
конечные автоматы для сопоставления имён, порядок обработки файлов,
инварианты многопоточности, управление ресурсами и модель ошибок.
\end{abstract}
\tableofcontents
\newpage

\section{Назначение и техническая среда}

Утилита формирует детерминированный манифест пар «относительный путь ---
криптографическая хэш-сумма содержимого». Корневой каталог задаётся
позиционным аргументом или по умолчанию равен текущему. Результат выводится в
\texttt{stdout}, сообщения об ошибках --- в \texttt{stderr}.

Собственная производственная часть использует C17 (ISO/IEC 9899:2018)
и POSIX-интерфейсы. Сторонние компоненты: Concurrency Kit 0.7.2,
Mbed TLS 3.6.7 LTS, GoogleTest 1.18.0 (только тесты). Сборка должна
работать без сетевых загрузок на macOS и Linux с поддерживаемыми
архитектурами процессора.

\subsection{Границы модулей}

\begin{longtable}{P{0.20\textwidth}P{0.71\textwidth}}
\toprule
Модуль & Основная ответственность \\
\midrule
\endhead
\texttt{matcher} & Компиляция wildcard-шаблонов, NFA, DFA, минимизация и сопоставление полного базового имени. \\
\texttt{queue} & Многопоточная MPMC-очередь заданий, жизненный цикл, согласованная обработка заполнения. \\
\texttt{catalog} & Локальные наборы путей, перенос владения, сортировка, состояния обработки и ожидание результата. \\
\texttt{hash} & Потоковое SHA-256, SHA-384, SHA-512 и перевод digest в HEX. \\
\texttt{hash\_pool} & Распределение индексов каталога между хэширующими рабочими потоками. \\
\texttt{walker} & Рекурсивное обнаружение файлов и каталогов с применением фильтров. \\
\texttt{manifest} & Ожидание упорядоченных результатов, экранирование имён и запись в stdout. \\
\texttt{app} & Обработка CLI, объединение библиотек, итоговый статус процесса и сигналов ОС. \\
\bottomrule
\end{longtable}

\begin{center}
\begin{tikzpicture}[
    box/.style={draw,rounded corners,minimum width=33mm,minimum height=10mm,align=center},
    >=Stealth,node distance=12mm and 14mm]
\node[box] (walk) {walker + matcher};
\node[box,right=of walk] (catalog) {catalog + qsort};
\node[box,below=of catalog] (hash) {hash\_pool + hash};
\node[box,left=of hash] (manifest) {manifest → stdout};
\draw[->] (walk) -- node[above]{пути} (catalog);
\draw[->] (catalog) -- node[right]{индексы} (hash);
\draw[->] (hash) -- node[below]{результаты} (manifest);
\end{tikzpicture}
\end{center}

Две рабочие стадии выполняются последовательно: сначала все пути обнаруживаются
и глобально сортируются, затем независимо вычисляются хэши. Печатать
сортированные готовые результаты допускается до завершения всех вычислений.

\section{Обнаружение и фильтрация}

\subsection{Обход каталогов}

Каждая рабочая единица содержит каталог для обследования. Walker-поток берёт
задание, открывает свой дескриптор каталога, перебирает элементы, фильтрует
базовые имена, добавляет подкаталоги в очередь, а имена файлов в свой локальный
массив. После завершения потоков локальные массивы объединяются и сортируются.

Счётчик незавершённых задач увеличивается \emph{перед} публикацией
каждого нового каталога и уменьшается \emph{после} окончания его обработки.
Нулевой счётчик одновременно означает, что нет ожидающих и активных заданий.
Само наличие временно пустой очереди недостаточно для завершения потоков.

\subsection{Ссылки и изменяемая файловая система}

Ссылка на обычный файл включается под именем ссылки, её содержимое получается
по целевому файлу. Ссылка на каталог не обходится; для неё выдаётся
диагностика. Повреждённые ссылки и ошибки чтения фиксируются; доступная часть
дерева продолжает обрабатываться.

Проверка типа объекта и последующее открытие могут наблюдать разные состояния
при конкурентном изменении файловой системы. Применяются \texttt{lstat},
\texttt{fstatat}, \texttt{fstat} и контроль результата открытия, но
\texttt{mangen} не создаёт атомарный снимок файловой системы.

\subsection{Matcher: модель wildcard-шаблона}

Модуль \texttt{matcher} проверяет \emph{имя файла целиком}.

Шаблон заранее преобразуется в конечный автомат.

Единица сопоставления --- \emph{расширенный
графемный кластер Unicode}, то есть воспринимаемый пользователем символ,
который может занимать несколько кодовых точек и байтов UTF-8. Например,
буква вместе с комбинируемым акцентом считается одним кластером.
Разбиение выполняет \texttt{utf8proc}; корректные участки UTF-8
нормализуются в форму NFC. Некорректные байты UTF-8 сохраняются и при
разбиении образуют отдельные единицы. Сопоставление чувствительно к регистру.

\begin{longtable}{P{0.21\textwidth}P{0.70\textwidth}}
\toprule
Элемент шаблона & Действие \\
\midrule
\endhead
Литерал & Совпадает с соответствующим графемным кластером после NFC. \\
\texttt{.} & Совпадает ровно с одним графемным кластером. \\
\texttt{*} & Совпадает с произвольным числом кластеров, включая ноль. \\
\texttt{\textbackslash} & Экранирует следующий графемный кластер, превращая его в литерал. \\
\bottomrule
\end{longtable}

Несколько шаблонов \texttt{-e} применяются независимо: имя исключается,
если хотя бы один скомпилированный шаблон совпал с ним. Сопоставляется
только \texttt{basename}, а не полный путь. Одинаковые соседние звёздочки
сворачиваются в одну при разборе; завершающая обратная косая черта
считается некорректным шаблоном.

\subsubsection{Математическая семантика}

Обозначим через $\Sigma$ алфавит нормализованных графемных кластеров,
дополненный отдельными некорректными байтами UTF-8. Через $\varepsilon$
обозначим пустую последовательность, а через $L(p)$ --- язык,
описываемый шаблоном $p$. Для литерала $g\in\Sigma$ выполняется:
\begin{align*}
 L(\varepsilon)&=\{\varepsilon\}, &
 L(g)&=\{g\}, \\
 L(\texttt{.})&=\Sigma, &
 L(\texttt{*})&=\Sigma^{*}, \\
 L(pq)&=L(p)\cdot L(q).
\end{align*}
Здесь $\cdot$ означает конкатенацию языков; для экранированного символа
применяется правило литерала. Если после NFC входное имя разбито на
$x=g_1g_2\ldots g_m$, то результат равен
\[
 \operatorname{match}(p,x)\iff x\in L(p).
\]
Поэтому шаблон \texttt{a*b} принимает \texttt{ab}, \texttt{axb} и
\texttt{axxxb}, но отвергает \texttt{aba}; шаблон \texttt{a.b}
принимает между \texttt{a} и \texttt{b} ровно один кластер.
Это полное совпадение, а не поиск подстроки.

\subsubsection{Разбор и построение NFA}

Компиляция начинается с токенизации
\texttt{MATCHER\_TOKEN\_LITERAL}, \texttt{ANY\_ONE},
\texttt{ANY\_MANY}. Непрерывные литералы нормализуются в NFC и
разбиваются на кластеры. Далее \texttt{matcher\_nfa\_build}
создаёт недетерминированный конечный автомат
$N=(Q,\Sigma,\delta,q_0,F)$ без явных $\varepsilon$-переходов.

Пусть $r$ --- число токенов, потребляющих ровно один кластер
(литерал или точка). Тогда
$Q=\{q_0,\ldots,q_r\}$, $F=\{q_r\}$.
После каждого потребляющего токена текущий номер состояния увеличивается:
\begin{align*}
 \delta(q_i,g)&\supseteq\{q_{i+1}\}
 &&\text{для литерала }g,\\
 \delta(q_i,c)&\supseteq\{q_{i+1}\}
 &&\text{для точки, любого }c\in\Sigma,\\
 \delta(q_i,c)&\supseteq\{q_i\}
 &&\text{для звёздочки, любого }c\in\Sigma.
\end{align*}
Звёздочка не добавляет состояние. Её петля позволяет поглощать
произвольное число кластеров, а переход следующего токена доступен
из того же состояния, поэтому нулевое число кластеров также допустимо.

\begin{center}
\begin{tikzpicture}[>=Stealth,node distance=25mm,
 state/.style={draw,circle,minimum size=11mm}]
\node[state] (q0) {$q_0$};
\node[state,right=of q0] (q1) {$q_1$};
\node[state,right=of q1,double] (q2) {$q_2$};
\draw[->] ([xshift=-11mm]q0.west) -- (q0);
\draw[->] (q0) -- node[above]{a} (q1);
\draw[->] (q1) to[out=125,in=55,looseness=4] node[above]{$\Sigma$} (q1);
\draw[->] (q1) -- node[above]{b} (q2);
\end{tikzpicture}
\end{center}
\begin{center}\small NFA для шаблона \texttt{a*b}.\end{center}

\subsubsection{Детерминизация и алфавитные классы}

Алгоритм подмножеств преобразует NFA в DFA. Состояние DFA представляет
множество активных состояний NFA, изначально $S_0=\{q_0\}$; после
прочтения кластера $c$ вычисляется
\[
 \Delta(S,c)=\bigcup_{q\in S}\delta(q,c),\qquad
 \operatorname{accept}(S)\iff q_r\in S.
\]
Число возможных подмножеств в общем случае достигает $2^{r+1}$,
но компилятор прекращает построение при превышении лимита состояний.
Множества реализованы битовыми векторами \texttt{uint64\_t}; переходы
по точке и петлям звёздочки рассчитываются битовыми масками и сдвигами,
а одинаковые множества переиспользуются через хэш-таблицу.

Поскольку конкретный литерал важен для перехода, а все остальные
кластеры одинаковы по отношению к литералам шаблона, используется
сжатый алфавит. Пусть $g_1,\ldots,g_k$ --- различные литеральные
кластеры шаблона. Функция класса
\[
 \kappa(c)=
 \begin{cases}
  i,&c=g_i,\quad 1\le i\le k,\\
  0,&c\notin\{g_1,\ldots,g_k\}
 \end{cases}
\]
сводит потенциально неограниченный алфавит к $k+1$ классам.
Литералы хранятся в отсортированном массиве; для выбора класса
используется двоичный поиск. Переходы DFA хранятся в плотной
таблице; индекс элемента вычисляется по формуле
\[
  i=s\cdot k_{\mathrm{classes}}+c,
\]
где $s$ --- текущее состояние, $c$ --- класс входного кластера.

\subsubsection{Минимизация DFA}

\texttt{matcher\_min\_dfa\_build} удаляет различия между состояниями,
которые принимают одинаковые множества продолжений. Сначала состояния
разбиваются на принимающие и непринимающие. Далее алгоритм уточнения
разбиения разделяет блоки по обратным переходам для каждого класса
алфавита, пока разделение более невозможно. Для быстрого поиска
предшественников построена структура CSR --- массив смещений и смежный
массив исходных состояний. Равные по поведению состояния заменяются
одним состоянием минимального DFA; его начальное состояние получает
индекс ноль.

\subsubsection{Сопоставление и оценка ресурсов}

Функция \texttt{mangen\_matcher\_compile} выполняет четыре этапа:
разбор шаблона, построение NFA, детерминизацию и минимизацию DFA.
Временные представления освобождаются, а объект
\texttt{mangen\_matcher\_t} владеет только минимальным DFA.
Для нового имени \texttt{mangen\_matcher\_match} нормализует вход,
выделяет кластеры, определяет класс $\kappa$ и делает один переход DFA
на кластер. Для ASCII-строк предусмотрен путь без выделения буфера NFC.

При $m$ кластерах входа и $k$ различных литералах стоимость
прохода после нормализации --- $O(m\log(k+1))$ сравнений классов
и $O(m)$ табличных переходов; таблица переходов занимает
$O(s(k+1))$ элементов, где $s$ --- число состояний минимального DFA.
Компиляция затратнее сопоставления, поэтому один объект выгодно
переиспользовать для множества имён. Реализация ограничивает входной
шаблон 1024 байтами и промежуточный DFA 4096 состояниями;
превышение лимита возвращает \texttt{MANGEN\_MATCHER\_LIMIT\_EXCEEDED}.

API различает отсутствие совпадения и внутреннюю ошибку через
\texttt{mangen\_matcher\_match}; сокращённый вызов
\texttt{mangen\_matcher\_matches} возвращает \texttt{false}
в обоих случаях. Скомпилированный объект при проверке не изменяется;
его освобождение через \texttt{mangen\_matcher\_free} нельзя выполнять
одновременно с активными проверками по тому же указателю.

\section{Очередь и многопоточность}

\subsection{MPMC}

\texttt{ck\_ring} используется для обмена заданиями между несколькими
производителями и потребителями. Состояния «закрыто» и «отменено» отличаются:
закрытие прекращает приём новых задач, но допускает выгрузку ожидающих;
отмена требует остановки и соединения всех потоков перед освобождением памяти.

Функция создания очереди задаёт минимальную ёмкость; она может
округляться согласно ограничениям \texttt{ck\_ring} и остаётся
постоянной на протяжении жизни объекта. Операция \texttt{try\_push}
не выделяет память и при заполнении возвращает статус
\texttt{MANGEN\_QUEUE\_FULL}; элемент остаётся у вызывающего кода.
Успешная запись передаёт владение очереди, успешное извлечение ---
потребителю.

Освобождение очереди допускается только после завершения всех потоков,
которые к ней обращаются.

\subsection{Условия отсутствия гонок}

\begin{enumerate}
\item Каждый каталог обрабатывается ровно одним владельцем задания.
\item Каждый индекс хэширования выдаётся не более одному worker-потоку.
\item Хэш и ошибка публикуются до перехода записи в терминальное состояние.
\item Поток печати ждёт результата по условной переменной;
активное ожидание исключено.
\item Освобождение каталога и очереди разрешено только после завершения
пользующихся ими потоков.
\end{enumerate}

\section{Каталог и вычисление SHA-2}

\subsection{Владение ресурсами}

После добавления пути локальный batch владеет строкой. При успешном
объединении владение получает каталог; он отвечает за освобождение строк.
Указатель на запись из каталога является заимствованным и действителен до
удаления каталога. Структуры закрыты непрозрачными C-типами.

\subsection{Состояние записи}

\begin{center}
\begin{tikzpicture}[>=Stealth,node distance=10mm,
  st/.style={draw,rounded corners,minimum width=28mm,minimum height=9mm}]
\node[st] (p) {PENDING};
\node[st,right=of p] (w) {PROCESSING};
\node[st,right=of w,yshift=6mm] (s) {SUCCESS};
\node[st,right=of w,yshift=-6mm] (e) {ERROR};
\draw[->] (p) -- (w);
\draw[->] (w) -- (s);
\draw[->] (w) -- (e);
\end{tikzpicture}
\end{center}

\texttt{SUCCESS} содержит валидные бинарные байты digest длины 32, 48 или 64,
\texttt{ERROR} содержит причину отказа. Ни один поток не должен печатать
незавершённый результат или использовать фиктивный хэш при ошибке.

\subsection{Потоковое хэширование}

Mbed TLS получает содержимое файла блоками фиксированного размера. Операции
\texttt{open/read/close} проверяются, размер файла не определяет размер
оперативной памяти хэширования. SHA-256 возвращает 64 HEX-символа,
SHA-384 --- 96, SHA-512 --- 128; буквенные символы A--F заглавные.
Стандарт хэш-функций --- NIST FIPS 180-4.

\section{Формат манифеста и ошибки}

\begin{lstlisting}
relative/path.txt : A3F1...<FULL_HEX_DIGEST>
\end{lstlisting}

Разделитель --- ровно \texttt{~:~} с пробелом с обеих сторон. В пути
экранируются обратная косая черта, LF и CR. Сортировка должна быть глобальной,
лексикографической по байтам относительного пути. Сообщения диагностики идут
в \texttt{stderr} независимо от \texttt{stdout}. Статусы процесса:
\texttt{0} --- успешный проход, \texttt{1} --- ошибка выполнения или
частичный результат, \texttt{2} --- ошибка CLI.

\section{Сложность и ограничения}

Пусть $F$ --- количество обнаруженных файлов, $D$ --- количество каталогов,
$S$ --- совокупный размер прочитанного содержимого. Обнаружение требует
$O(F+D)$ посещений, сортировка --- $O(F\log F)$ сравнений путей,
хэширование --- $O(S)$ работы по обработке байтов. В памяти остаётся
$O(F)$ записей и путей, рабочая очередь и ограниченные буферы потоков.
Пиковое потребление зависит от длин путей и числа очередных задач.

\section{Проверяемая переносимость}

На macOS и Linux нужны POSIX threads, атомарные операции C17 и корректная
сборка Concurrency Kit для данной архитектуры. В CMake следует использовать
\texttt{find\_package(Threads REQUIRED)} и цель \texttt{Threads::Threads}.
Для \texttt{realpath(path, NULL)} на glibc при строгом C17 нужны подходящие
feature-test macros (например, \texttt{\_POSIX\_C\_SOURCE=200809L}),
если данный вызов действительно используется.

Подтверждённые пользовательским журналом результаты macOS и проверки Linux
разграничены в документе \texttt{verification.pdf}.

\clearpage
\addcontentsline{toc}{section}{Список литературы}
\begingroup
\footnotesize
\setlength{\parskip}{0pt}
\begin{thebibliography}{9}
\bibitem{c17} ISO/IEC 9899:2018. Programming languages --- C.
\bibitem{posix} The Open Group. POSIX Base Specifications. \url{https://pubs.opengroup.org/onlinepubs/9799919799/}
\bibitem{sha} NIST. FIPS 180-4. Secure Hash Standard. \url{https://csrc.nist.gov/pubs/fips/180-4/upd1/final}
\bibitem{ck} Concurrency Kit. \url{https://github.com/concurrencykit/ck}
\bibitem{tls} Mbed TLS. \url{https://github.com/Mbed-TLS/mbedtls}
\bibitem{cmake} CMake. FindThreads. \url{https://cmake.org/cmake/help/latest/module/FindThreads.html}
\end{thebibliography}
\endgroup
\end{document}
